% Nouveau EGA V (cote 25): archive cote 25, site page 40 (PDF page 41).
% typescript. Draft compared: none.
% AI-transcribed by Claude Opus 5.5 (claude-opus-5-5, adaptive thinking, effort high), Message Batches API, batch ega25b (msgbatch_01A7n5YuXubUPNK1zsepi8RG), 2026-11-13; not yet checked by eye.
\scanpage{25/40}
\mspage{\hand{36}}

\diagram{\begin{tikzcd}
 & \unsure{0} & & \hand{0} & & \\
 & \underline{\Omega}^1_{X/Y} \otimes \underline{O}_Z(1,0) \arrow[u] & & \hand{\unsure{\mathcal{T}}}\check{}_{P/S}\,\unsure{(\otimes\,\underline{O}_Z(0,-1)} \arrow[u] & & \\
0 \arrow[r] & \underline{\Omega}^1_{P/S} \otimes \underline{O}_Z(1,0) \arrow[r] \arrow[u] & E\otimes\underline{O}_Z \arrow[r] \arrow[u] & \underline{O}_Z(1,0) \arrow[r] & 0 \\
 & \hand{\unsure{\check{\mathcal{N}}}}_{X/P} \otimes \underline{O}_Z(1,0) \arrow[u] & \underline{O}_Z(0,-1) \arrow[u] \arrow[ul, dashed, "\hand{\beta}"'] \arrow[ur, "\hand{\alpha}"'] & & & \\
 & 0 \arrow[u] & \hand{0} \arrow[u] & & &
\end{tikzcd}}

où la première colonne est déduite de (a) en tensorisant par $\underline{O}_Z(1,0)$,
la ligne est déduite de  (b) en tensorisant par $\underline{O}_Z$  , et la colonne 2
est déduite  de la \inserted{transposée de la} suite analogue \xout{xxxxxxx} (b$\hand{^\vee}$ ) relative à $\check{P}$  (obtenue
en remplaçant $\underline{E}$ par $\hand{\check{\underline{E}}}$ )  en tensorisant par $\underline{O}_Z$ . Par définition même
de $\underline{Y}$, Z se trouve au dessus de $\underline{Y}$  si et seulement si le morphisme
composé \hand{$\alpha$} du diagramme est nul, i.e.  si on peut trouver une  factorisa
tion \hand{$\beta$}: $\underline{O}_Z(0,-1) \rightarrow \underline{\Omega}^1_{P/S}\otimes \underline{O}_Z(1,0)$ .  S'il en est ainsi, on  peut
considérer  son composé avec  $\underline{\Omega}^1_{P/S} \otimes \underline{O}_Z(1,0) \rightarrow \underline{\Omega}^1_{X/Y} \otimes \underline{O}_Z(1,0)$,
qu'on peut aussi interpréter comme une section de $\underline{\Omega}^1_{X/Y}\otimes \underline{O}_Z(1,1)$.
Je dis que c'est précisément la section \inserted{$\unsure{\nu|\mu}$} qu'on avait envisagé plus haut.
(La vérification doit être essentiellement mécanique).  Elle est
nulle  si et seulement si  Z est au dessus de V(\hand{$\unsure{\nu|\mu}$})  (par définition
même de V(\hand{$\unsure{\nu|\mu}$}) !) , mais cela signifie aussi  que \hand{$\beta$} se factorise par
$\hand{\unsure{\check{\mathcal{N}}}}_{X/P} \otimes \underline{O}_Z(1,0)$, i.e. que  le sous-Module $\underline{O}_Z(0,-1)$ de $E\otimes \underline{O}_Z$ est contenu
dans le sous-Module $\hand{\unsure{\mathcal{N}}}_{X/P} \otimes \underline{O}_Z(1,0)$,  ce qui évidemment signifie aussi
que  Z se trouve au dessus du sous-préschéma  P( $\hand{\unsure{\mathcal{N}}}_{X/P}$ (1)) de P($\hand{\check{\underline{E}}}_X$) , et
achève la démonstration de ce que nous avions annoncé.

Dès avant ce savant exercice de syntaxe (qui m'a couté d'ailleurs

\margin{\hand{40}}
